\documentclass[12pt,a4paper]{report}
\usepackage{lmodern} 
\usepackage[utf8]{inputenc}
\usepackage[T1]{fontenc}
\usepackage[spanish]{babel}
\usepackage{hyperref}
\usepackage{microtype} 
\usepackage{xurl}      

\begin{document}

\title{Cognia}
\author{Pavel Bespakhotnyi \\ Alla Kolchanova}
\date{\today}
\maketitle

\tableofcontents

\chapter{Introducción, problema y objetivos}

\section{Introducción}

Cognia es una plataforma tecnológica multiplataforma diseñada para
transformar la gestión, la organización pedagógica y la comunicación en
el sector de las clases particulares y las tutorías independientes. La
solución unifica en un único ecosistema coordinado a los tres actores
principales del proceso educativo: el profesor particular, el alumno y
sus progenitores o tutores legales.

El nombre del proyecto, Cognia, proviene del latín \textit{cognitio}
(conocimiento, aprendizaje, conciencia). El eslogan elegido es
\textit{«Clases estructuradas, progreso visible»}, que resume la
propuesta de valor: ofrecer al profesor una herramienta profesional de
gestión y al alumno y su familia una visibilidad clara del progreso
académico.

\section{Problema o necesidad detectada}

En el modelo tradicional de clases particulares, el profesor autónomo
gestiona de forma simultánea entre 10 y 25 alumnos. La ausencia de un
software especializado obliga al docente a recurrir a herramientas
genéricas e inconexas: hojas de cálculo aisladas (Excel o Google Sheets)
para los calendarios, chats de WhatsApp para avisos y recepción de
deberes, e impresos o notas en libreta para el control financiero.

Esta fragmentación operativa genera tres problemas concretos:

\begin{itemize}
    \item \textbf{Para el profesor:} dedica entre 5 y 10 horas semanales
    a tareas administrativas (cuadrar horarios, controlar cobros, enviar
    tareas, redactar informes), tiempo que resta a la preparación de sus
    contenidos didácticos.
    \item \textbf{Para el alumno:} pierde el seguimiento de sus deberes
    entre cientos de mensajes, no tiene visibilidad de sus entregas
    pendientes ni de su evolución, y pierde la motivación.
    \item \textbf{Para el padre o tutor legal:} sufre una opacidad total
    sobre el rendimiento real de su hijo y se entera de los impagos o de
    las deudas acumuladas con meses de retraso.
\end{itemize}

\subsection*{Situaciones reales de conflicto}

\begin{itemize}
    \item \textbf{Caso A (conflicto financiero):} el padre de un alumno
    recibe una notificación verbal de impago acumulado de dos meses
    (240~€) porque el profesor olvidó anotar tres sesiones dictadas
    durante días festivos y no emitió un recordatorio formal a tiempo.
    \item \textbf{Caso B (pérdida de material):} una estudiante de 3.º de
    ESO no entrega su trabajo de preparación de examen porque el enlace
    con la ficha PDF se perdió en un hilo de WhatsApp con más de 200
    mensajes cruzados.
    \item \textbf{Caso C (saturación administrativa):} un profesor
    invierte la noche del domingo (más de 3 horas) redactando
    manualmente 15 correos individuales a los padres para informar sobre
    las notas del examen semanal.
\end{itemize}

\subsection*{Limitaciones de las soluciones actuales}

Las herramientas existentes fallan al aplicarse al contexto del tutor
privado independiente:

\begin{itemize}
    \item \textbf{Aplicaciones de mensajería (WhatsApp, Telegram):}
    mezclan la vida personal con la profesional. La información académica
    se diluye en hilos informales, los archivos adjuntos caducan y no
    existe estructura de seguimiento.
    \item \textbf{Hojas de cálculo (Excel, Sheets):} requieren
    actualización manual constante, sufren errores humanos y no ofrecen
    una interfaz interactiva ni notificaciones automáticas.
    \item \textbf{Plataformas escolares institucionales (Google
    Classroom, Moodle, Alexia):} están diseñadas para colegios e
    institutos complejos. Son hipertrofiadas, requieren configuraciones
    de dominio institucional y no contemplan la relación comercial ni los
    cobros del tutor privado.
\end{itemize}

\section{Usuario o cliente objetivo}

Cognia define tres perfiles de usuario perfectamente segregados en roles,
competencias técnicas y necesidades operativas:

\begin{itemize}
    \item \textbf{Profesor particular (ROLE\_TUTOR):} autónomo que
    gestiona entre 5 y 20 alumnos. Edad: 25--50 años. Nivel técnico:
    básico--medio. Utiliza un panel web integral desde su ordenador
    portátil al finalizar su jornada académica.
    \item \textbf{Alumno (ROLE\_STUDENT):} estudiante de entre 10 y 18
    años. Nivel técnico: alto (nativo digital). Utiliza una aplicación
    móvil Android desde casa, en el centro educativo o en movimiento.
    \item \textbf{Padre o tutor legal (ROLE\_PARENT):} adulto de entre 35
    y 55 años. Nivel técnico: básico. Consulta el historial académico de
    su hijo desde su dispositivo móvil o desde el correo electrónico.
\end{itemize}

\subsection*{User Personas representativas}

\begin{itemize}
    \item \textbf{Carlos Ruiz (profesor particular):} 34 años,
    licenciado en Matemáticas, tutor privado de 18 alumnos de secundaria.
    Necesita centralizar su agenda, reducir el tiempo dedicado a cobros y
    enviar informes profesionales a los padres sin perder horas en
    correos manuales.
    \item \textbf{Lucía Gómez (alumna de 3.º de ESO):} 15 años,
    estudiante que requiere refuerzo en Física y Química. Necesita ver
    claramente sus deberes pendientes, descargar las fichas de ejercicios
    en su móvil y sentirse motivada al ver su progreso.
    \item \textbf{Roberto Gómez (padre de Lucía):} 48 años, ingeniero
    industrial con jornada laboral intensa. Necesita tener visibilidad
    del avance académico de su hija y recibir avisos claros sobre los
    pagos mensuales sin tener que preguntar constantemente.
\end{itemize}

\section{Objetivos generales del proyecto}

Los objetivos de Cognia se han formulado bajo la metodología SMART
(específicos, medibles, alcanzables, relevantes y temporales),
garantizando su verificación cuantitativa al finalizar el proyecto:

\begin{enumerate}
    \item \textbf{Optimizar el tiempo administrativo del profesor.}
    Métrica: tiempo total por sesión (registro de clase, tarea e
    informe). Criterio de éxito: inferior a 2 minutos, lo que supone una
    reducción del 80\,\% respecto al método manual.

    \item \textbf{Facilitar el acceso y la descarga de tareas en la
    aplicación Android.} Métrica: pruebas de rendimiento con Android
    Profiler. Criterio de éxito: visualización de tareas y descarga del
    PDF en menos de 1,5 segundos.

    \item \textbf{Automatizar las comunicaciones con las familias.}
    Métrica: notificaciones y correos enviados. Criterio de éxito: el
    100\,\% de los informes se envían en menos de 30 segundos tras el
    registro de la clase.

    \item \textbf{Garantizar la disponibilidad y el rendimiento de la
    API REST.} Métrica: pruebas de carga con Postman. Criterio de éxito:
    disponibilidad superior al 99\,\% y tiempo de respuesta inferior a
    300~ms en el 95\,\% de las peticiones.

    \item \textbf{Asegurar el funcionamiento sincronizado del sistema
    en las plataformas web y móvil.} Métrica: pruebas de integración
    cruzada. Criterio de éxito: los tres canales de acceso (web, Android
    y correo electrónico) operan de forma sincronizada sin errores.
\end{enumerate}

\section{Justificación del proyecto}

Cognia resuelve un problema real que afecta a las tres partes
involucradas en el proceso educativo de las clases particulares. Frente a
las alternativas existentes, el proyecto aporta las siguientes ventajas:

\begin{itemize}
    \item Centraliza en una única plataforma la gestión administrativa,
    pedagógica y comunicativa del profesor autónomo.
    \item Ofrece al alumno una experiencia móvil clara y motivadora,
    con acceso directo a sus tareas, fechas límite y progreso.
    \item Proporciona a las familias transparencia total sobre el
    rendimiento académico y contable de sus hijos.
    \item Se ajusta perfectamente al plazo de 13 \textit{sprints}
    establecido para el proyecto intermodular del ciclo DAM2.
    \item No requiere infraestructura en tiempo real ni APIs externas
    complejas, lo que garantiza su viabilidad técnica y económica.
\end{itemize}

\chapter{Análisis de Requisitos y Control de Acceso}

\section{Introducción al análisis funcional}

Una vez definidos el problema, los usuarios y los objetivos generales en el primer capítulo, el siguiente paso es traducir esas necesidades en un modelo funcional concreto. Este capítulo detalla cómo se estructuran los roles dentro del sistema, las reglas de negocio para garantizar la seguridad y el listado formal de requisitos que guiarán el desarrollo técnico de Cognia.

\section{Definición de Roles y Gestión de Identidades}

Para garantizar la privacidad de los datos (especialmente los de menores de edad) y evitar el caos operativo, se ha establecido una arquitectura de accesos estrictamente controlada:

\begin{itemize}
    \item \textbf{Roles del sistema:} La plataforma define tres roles principales: \texttt{TUTOR}, \texttt{STUDENT} y \texttt{PARENT}.
    \item \textbf{Registro propio de los tres roles:} Cada usuario crea su propia cuenta según su rol: el profesor en la aplicación web, el alumno y el padre en la aplicación Android. El profesor no crea cuentas de alumnos ni de padres, y ningún rol puede crear cuentas de otro rol.
    \item \textbf{Vinculación pendiente:} Las relaciones entre alumno y profesor, y entre alumno y padre, se definirán antes de implementar la API (código de invitación, enlace por correo u otro).
    \item \textbf{Redirección automática basada en el Backend:} Tras el inicio de sesión, el backend (Spring Boot) validará el token JWT y determinará el rol del usuario. El sistema lo redirigirá automáticamente a su interfaz correspondiente (Web para el profesor, Android para el alumno o el padre), imposibilitando el acceso cruzado a datos que no le corresponden.
\end{itemize}

\section{Requisitos Funcionales}

Los requisitos funcionales describen las acciones que el sistema debe permitir. Se han clasificado según su prioridad utilizando el método MoSCoW (\textit{Must have}, \textit{Should have}, \textit{Could have}).

\subsection{Módulo Web (Profesor particular)}
\begin{itemize}
    \item \textbf{RF-01 (Must have):} Registro propio del profesor en la aplicación web y autenticación segura mediante token JWT.
    \item \textbf{RF-02 (Must have):} Gestión CRUD (Crear, Leer, Actualizar, Borrar) de las fichas de los alumnos y sus datos de contacto.
    \item \textbf{RF-03 (Must have):} Calendario interactivo para agendar clases, con posibilidad de cancelar o reprogramar.
    \item \textbf{RF-04 (Must have):} Registro y control del estado financiero de cada clase (Pagado / Pendiente / Deuda).
    \item \textbf{RF-05 (Must have):} Subida de material didáctico y tareas en formato PDF para cada sesión.
    \item \textbf{RF-06 (Should have):} Generador de informes de sesión que se envía automáticamente al correo del padre.
    \item \textbf{RF-07 (Could have):} Panel de estadísticas mensuales (ingresos, clases impartidas, ratio de asistencia).
\end{itemize}

\subsection{Módulo Android (Alumno y Padre)}
\begin{itemize}
    \item \textbf{RF-08 (Must have):} Inicio de sesión unificado. El sistema detecta el rol y muestra la interfaz correcta sin necesidad de seleccionar un perfil manualmente.
    \item \textbf{RF-09 (Must have - Alumno):} Visualización de la lista de tareas pendientes y descarga de archivos PDF adjuntos.
    \item \textbf{RF-10 (Should have - Alumno):} Historial de calificaciones y visualización del progreso académico.
    \item \textbf{RF-11 (Should have - Padre):} Consulta del historial de clases impartidas a su hijo y acceso a los informes del profesor.
    \item \textbf{RF-12 (Could have - Padre):} Recepción de notificaciones push sobre la publicación de nuevas notas o avisos de pago.
    \item \textbf{RF-13 (Could have - Alumno):} Sistema de gamificación (acumulación de puntos por tareas entregadas a tiempo).
    \item \textbf{RF-14 (Must have - Padre):} Registro propio del padre o tutor legal en la aplicación Android, sin intervención del profesor.
    \item \textbf{RF-15 (Must have - Alumno):} Registro propio del alumno en la aplicación Android, sin intervención del profesor.
\end{itemize}

\section{Requisitos No Funcionales}

Los requisitos no funcionales definen los estándares de calidad y rendimiento que debe cumplir la plataforma:

\begin{itemize}
    \item \textbf{Rendimiento:} La API REST debe procesar y responder a las peticiones en menos de 300 milisegundos en el 95\% de los casos. El tiempo de carga inicial del portal Web no debe superar los 2 segundos.
    \item \textbf{Seguridad:} 
    \begin{itemize}
        \item Uso de HTTPS (TLS 1.3) para todo el tráfico entre los clientes y el servidor.
        \item Cifrado de contraseñas en la base de datos mediante el algoritmo BCrypt.
        \item Protección de los endpoints de la API mediante tokens JWT (con tiempo de expiración definido).
    \end{itemize}
    \item \textbf{Usabilidad:} Interfaz en español, diseño responsive adaptado a los estándares de accesibilidad visual WCAG 2.1 (nivel AA).
    \item \textbf{Compatibilidad:} Soporte garantizado para navegadores modernos (Chrome, Firefox, Edge) y dispositivos móviles con Android 8.0 o superior.
    \item \textbf{Disponibilidad:} El servidor VPS debe garantizar una disponibilidad del 99\% durante el periodo lectivo de producción.
    \item \textbf{Mantenibilidad:} El código fuente debe estar documentado (JavaDoc/JSDoc) y versionado en GitHub bajo la metodología GitFlow.
\end{itemize}

\section{Análisis de Riesgos Inicial}

Durante la fase de planificación de requisitos, se han identificado los siguientes riesgos técnicos que podrían afectar al desarrollo:

\begin{enumerate}
    \item \textbf{Complejidad de la seguridad:} La configuración de Spring Security con múltiples roles y JWT puede alargar el Sprint de Backend. \textit{Mitigación:} Realizar pruebas tempranas con Postman.
    \item \textbf{Gestión de archivos PDF:} El peso de los archivos puede ralentizar la app Android. \textit{Mitigación:} Limitar los archivos a un máximo de 5 MB y usar Retrofit para descargas asíncronas.
    \item \textbf{Sincronización de la base de datos:} Errores en las relaciones entre tablas (ej. alumnos sin padres asignados). \textit{Mitigación:} Diseñar el modelo Entidad-Relación en 3FN antes de escribir código en Spring Boot.
\end{enumerate}

\chapter{Configuración técnica del proyecto y estructura del repositorio}

\section{Introducción}

Una vez definidos los requisitos funcionales y no funcionales en el capítulo anterior, el siguiente paso es establecer las bases técnicas sobre las que se construirá Cognia. Este capítulo documenta la configuración inicial del repositorio, la organización de carpetas, la instalación de las herramientas de desarrollo y la preparación de los entornos de trabajo tanto para el frontend web como para los futuros módulos de backend y móvil.

El objetivo de esta fase no es implementar funcionalidad de negocio, sino garantizar que cualquier miembro del equipo pueda clonar el repositorio, instalar las dependencias y ejecutar el proyecto en su máquina local con un único comando. Se han seguido principios de \textit{monorepositorio}, separación de responsabilidades (\textit{separation of concerns}) y automatización de tareas repetitivas.

\section{Estructura del repositorio}

\subsection{Enfoque de monorepositorio}

Se ha adoptado una estructura de \textit{monorepositorio}: un único repositorio de Git contiene todos los componentes del proyecto (documentación, backend, frontend web, aplicación Android y recursos compartidos). Esta decisión se justifica por:

\begin{itemize}
    \item \textbf{Visión unificada:} cualquier miembro del equipo accede a todo el código con un solo \texttt{git clone}.
    \item \textbf{Versionado coherente:} los cambios en la API y en el cliente web se revisan en el mismo \textit{pull request}.
    \item \textbf{Simplicidad de CI/CD:} una sola configuración de integración continua gestiona todos los módulos.
\end{itemize}

\subsection{Estructura de carpetas}

La estructura final del repositorio es la siguiente:

\begin{verbatim}
cognia/
|-- docs/               # Documentacion (BD, LaTeX, Figma, OpenAPI)
|-- backend/            # Spring Boot (Java) - REST API
|-- frontend-web/       # Vite + React + TypeScript
|-- mobile-android/     # Android (Java + Retrofit)
|-- shared/             # OpenAPI y design-tokens
|-- infra/              # docker-compose, nginx, CI/CD
|-- .gitignore
`-- README.md
\end{verbatim}

Cada carpeta tiene una responsabilidad claramente definida:

\begin{itemize}
    \item \path{docs/}: documentación LaTeX, esquema DBML, especificación OpenAPI y exportaciones de Figma.
    \item \path{backend/}: albergará el servidor Spring Boot con la API REST.
    \item \path{frontend-web/}: contiene la aplicación web desarrollada con Vite, React y TypeScript, destinada exclusivamente al rol \texttt{TUTOR}.
    \item \path{mobile-android/}: contendrá la aplicación Android unificada para alumnos y padres.
    \item \path{shared/}: recursos compartidos entre módulos, como los contratos OpenAPI y los tokens de diseño (\textit{design tokens}) exportados desde Figma.
    \item \path{infra/}: configuración de infraestructura: \path{docker-compose.yml}, configuración de Nginx y archivos de integración continua.
\end{itemize}

\subsection{Gestión de versiones y convenciones}

Se ha configurado un archivo \path{.gitignore} en la raíz del repositorio que excluye de forma selectiva los artefactos generados por cada tecnología:

\begin{itemize}
    \item \textbf{Node.js / Vite:} \path{node_modules/}, \path{dist/}, \path{storybook-static/}, \path{.vite/}.
    \item \textbf{Java / Gradle:} \path{build/}, \path{.gradle/}, \texttt{*.class}, \texttt{*.jar}.
    \item \textbf{Android:} \path{local.properties}, \texttt{*.apk}, \texttt{*.aab}, \texttt{*.keystore}.
    \item \textbf{Secretos:} \path{.env}, \path{.env.*} (excepto \path{.env.example}).
    \item \textbf{LaTeX:} archivos auxiliares (\path{.aux}, \path{.log}, \path{.toc}, \path{.synctex.gz}).
\end{itemize}

Además, se ha añadido un archivo \path{docs/latex/.gitignore} específico para excluir los archivos temporales generados por la compilación de LaTeX, manteniendo únicamente el código fuente (\path{main.tex}) en el repositorio.

\section{Preparación del frontend web}

\subsection{Inicialización con Vite, React y TypeScript}

La aplicación web se ha inicializado mediante Vite
con la plantilla \texttt{react-ts}: 

\begin{verbatim}
npm create vite@latest . -- --template react-ts
npm install
\end{verbatim}

Esta elección proporciona:

\begin{itemize}
    \item \textbf{React 18+}: biblioteca de interfaz de usuario basada en componentes.
    \item \textbf{TypeScript}: tipado estático que reduce errores en tiempo de compilación.
    \item \textbf{Vite}: servidor de desarrollo con recarga en caliente (\textit{HMR}) y empaquetado optimizado para producción.
\end{itemize}

El servidor de desarrollo se ejecuta en \url{http://localhost:5173} mediante el comando \texttt{npm run dev}.

\subsection{Configuración de ESLint y Prettier}

Para garantizar la calidad y consistencia del código se han configurado dos herramientas complementarias:

\begin{itemize}
    \item \textbf{ESLint} (\path{eslint.config.js}): análisis estático que
    detecta errores de programación, malas prácticas y problemas de tipado.
    Se ha utilizado
    la configuración plana (\textit{flat config}) introducida en ESLint 9, con los
    siguientes complementos: \path{@eslint/js}, \path{typescript-eslint},
    \path{eslint-plugin-react-hooks} y \path{eslint-plugin-react-refresh}.
    \item \textbf{Prettier} (\path{.prettierrc}): formateador de código que impone un estilo uniforme (comillas simples, punto y coma obligatorio, ancho máximo de 80 caracteres, indentación de 2 espacios).
\end{itemize}

Ambas herramientas se integran mediante \path{eslint-config-prettier}, que desactiva las reglas de ESLint que entrarían en conflicto con el formato de Prettier. Los comandos disponibles son:

\begin{verbatim}
npm run lint      # Ejecuta ESLint
npm run format    # Aplica Prettier
\end{verbatim}

\subsection{Instalación y configuración de Storybook}

Storybook se ha instalado con el comando:

\begin{verbatim}
npx storybook@latest init
\end{verbatim}

Storybook proporciona un entorno aislado para desarrollar y documentar componentes de interfaz de usuario. Se ha integrado el \texttt{addon-designs}, que permite vincular cada historia (\textit{story}) con su correspondiente diseño en Figma:

\begin{verbatim}
npx storybook@latest add @storybook/addon-designs
\end{verbatim}

La organización adoptada separa estrictamente los componentes de sus historias:

\begin{itemize}
    \item \path{src/components/}: componentes reutilizables, sin archivos \texttt{.stories}.
    \item \path{src/pages/}: pantallas completas de la aplicación.
    \item \path{src/stories/}: todas las historias de Storybook, que importan los componentes desde \path{src/components/}.
\end{itemize}

Esta separación evita que los archivos de documentación visual se mezclen con el código de producción. El servidor de Storybook se ejecuta en \url{http://localhost:6006} mediante el comando \texttt{npm run storybook}.

\section{Preparación de la infraestructura local}

\subsection{Contenedor de PostgreSQL con Docker Compose}

Para garantizar que todos los miembros del equipo dispongan de una base de datos idéntica sin necesidad de instalar PostgreSQL manualmente en el sistema operativo, se ha creado el archivo \path{infra/docker-compose.yml}. Este archivo define un único servicio:

\begin{verbatim}
services:
  postgres:
    image: postgres:16-alpine
    container_name: cognia-postgres
    environment:
      POSTGRES_DB: cognia
      POSTGRES_USER: cognia
      POSTGRES_PASSWORD: cognia_dev_password
    ports:
      - "5432:5432"
    volumes:
      - cognia_pgdata:/var/lib/postgresql/data
\end{verbatim}

La base de datos se inicia con el comando \texttt{docker compose up -d} y se detiene con \texttt{docker compose down}. El volumen \path{cognia_pgdata} garantiza la persistencia de los datos entre reinicios.

Es importante destacar que este contenedor únicamente proporciona un servidor PostgreSQL vacío: la estructura de tablas se creará automáticamente mediante las migraciones de Flyway incluidas en el backend (\path{backend/src/main/resources/db/migration/}).

\subsection{Documentación de la infraestructura}

Se ha creado el archivo \path{infra/README.md} en inglés que documenta:

\begin{itemize}
    \item El propósito de la carpeta \path{infra/}.
    \item Las credenciales de desarrollo (exclusivamente locales).
    \item Los comandos de Docker Compose (iniciar, detener, ver logs, abrir \texttt{psql}).
    \item La hoja de ruta para futuras incorporaciones (servicios \texttt{backend} y \texttt{nginx}).
\end{itemize}

\section{Documentación de la base de datos}

Se ha creado el archivo \path{docs/database/README.md} que documenta el esquema relacional de Cognia. Incluye:

\begin{itemize}
    \item Agrupación de las tablas en seis bloques funcionales: identidad, clases, indicadores del alumno, finanzas, comunicación y extensiones futuras.
    \item Descripción de las relaciones entre tablas (claves foráneas, cardinalidades).
    \item Convenciones del proyecto (UUID como clave primaria, \texttt{decimal} para importes monetarios, estados almacenados como \texttt{varchar}).
    \item Procedimiento para convertir el diseño DBML en migraciones Flyway.
\end{itemize}

Este documento complementa al archivo \path{schema_v5.0.dbml}, que constituye la fuente de verdad del modelo de datos.

\section{Resultados de la fase de configuración}

Al finalizar esta fase, el repositorio contiene:

\begin{itemize}
    \item Una estructura de carpetas escalable y documentada.
    \item Un frontend web funcional con React, TypeScript, ESLint, Prettier y Storybook.
    \item Una base de datos PostgreSQL reproducible localmente mediante Docker Compose.
    \item Un archivo \path{.gitignore} que protege al repositorio de artefactos y secretos.
    \item Documentación técnica completa en \path{infra/README.md} y \path{docs/database/README.md}.
\end{itemize}

Todos los cambios han sido versionados en la rama \texttt{main} mediante commits atómicos y descriptivos. El repositorio está preparado para que el siguiente \textit{sprint} se centre en el diseño en Figma de la interfaz (UI Kit y pantallas de autenticación) y en la implementación de las primeras pantallas del frontend.

\chapter{Desarrollo e implementación del Frontend}

\section{Introducción}

Una vez finalizada la configuración técnica del repositorio y del entorno de desarrollo (descrita en el capítulo anterior), el siguiente paso consistió en construir la interfaz de usuario de la aplicación web. Este capítulo documenta el proceso de transformación de los diseños de Figma en componentes funcionales de React, la organización del código para garantizar su escalabilidad y la implementación de la lógica de navegación entre las distintas pantallas de la plataforma.

El objetivo de esta fase es doble: por un lado, proporcionar una experiencia de usuario (UX) clara y profesional para el rol del profesor particular (\texttt{TUTOR}), y por otro, establecer una arquitectura de componentes reutilizables que facilite el futuro desarrollo del panel de control (Dashboard) y la integración con la API REST de Spring Boot.

\section{Arquitectura de enrutamiento y organización de páginas}

\subsection{Transición de un prototipo monolítico a una SPA}

En las primeras fases de prototipado, la totalidad de las vistas se gestionaba desde un único archivo \path{App.tsx} mediante el uso del hook \texttt{useState}. Aunque este enfoque es rápido para validar ideas, resulta inviable para un proyecto real, ya que provoca un archivo excesivamente largo y dificulta el mantenimiento.

Para resolver esta limitación, se ha implementado el enrutamiento del lado del cliente (\textit{Client-Side Routing}) utilizando la biblioteca \textbf{React Router}. Esto convierte a Cognia en una Aplicación de Página Única (\textit{Single Page Application}, SPA), donde las transiciones entre pantallas son instantáneas, sin recargar el navegador.

\subsection{Estructura de directorios escalable}

Para mantener el código limpio y predecible, se ha adoptado la siguiente estructura de carpetas dentro de \path{src/}:

\begin{itemize}
    \item \path{components/}: Contiene componentes globales y reutilizables que se repiten en varias páginas, como \texttt{Header.tsx} y \texttt{Footer.tsx}.
    \item \path{pages/}: Contiene los archivos que representan una pantalla completa de la aplicación. Cada archivo es responsable de ensamblar los componentes y gestionar su propia lógica local.
    \item \path{assets/}: Almacena los recursos estáticos (imágenes, logotipos, fotografías) exportados desde Figma.
    \item \path{stories/}: Reservado exclusivamente para las historias de Storybook.
\end{itemize}

Esta separación de responsabilidades (\textit{Separation of Concerns}) permite que, al añadir nuevas funcionalidades (como el calendario o la gestión de alumnos), el código no se mezcle y sea fácil localizar dónde realizar cambios.

\subsection{Implementación de las rutas}

El archivo principal \path{App.tsx} se ha refactorizado para actuar únicamente como un despachador de rutas. Utilizando los componentes \texttt{BrowserRouter} y \texttt{Routes}, se han definido las siguientes vistas:

\begin{itemize}
    \item \texttt{/} (o raíz): Muestra la página de inicio (\texttt{LandingPage}).
    \item \texttt{/login}: Muestra el formulario de inicio de sesión (\texttt{LoginPage}).
    \item \texttt{/register}: Muestra el formulario de registro con selección de roles (\texttt{RegisterPage}).
    \item \texttt{/forgot-password}: Muestra la pantalla de recuperación de contraseña (\texttt{ForgotPasswordPage}).
    \item \texttt{/email-verification}: Muestra la pantalla de confirmación tras el registro (\texttt{EmailVerificationPage}).
    \item \texttt{/about}: Muestra la información institucional del proyecto (\texttt{AboutPage}).
    \item \texttt{/dashboard}: Página de inicio de sesión exitosa (actualmente un marcador de posición para futuras iteraciones).
\end{itemize}

\section{Diseño de la interfaz y gestión de estilos}

\subsection{Actualización de la paleta de colores}

Durante la fase de diseño en Figma, se definió una paleta de colores corporativa que difiere de los colores por defecto generados por Vite. Para reflejar esta identidad visual, se actualizó el archivo \path{index.css} definiendo variables CSS (Custom Properties) en el selector \texttt{:root}. Los nuevos colores son:

\begin{itemize}
    \item \textbf{Fondo (Background):} \texttt{\#f6f8f7} (un tono gris muy claro, relajante para la vista).
    \item \textbf{Primario (Primary):} \texttt{\#28d745} (verde brillante, utilizado para llamadas a la acción).
    \item \textbf{Secundario (Secondary):} \texttt{\#789dfc} (azul suave, para textos destacados o botones secundarios).
    \item \textbf{Acento (Accent):} \texttt{\#b966e5} (violeta, reservado para etiquetas o detalles decorativos).
    \item \textbf{Texto (Text):} \texttt{\#000000} (negro puro para máxima legibilidad).
\end{itemize}

El uso de variables CSS (\texttt{var(--color-primary)}) permite que, si en el futuro se decide cambiar el verde por un azul, solo sea necesario modificar una línea en el archivo \path{index.css} para que toda la aplicación se actualice automáticamente.

\subsection{Separación de estilos: App.css vs estilos locales}

Para evitar conflictos de CSS, se ha aplicado una estrategia de separación:

\begin{itemize}
    \item \path{App.css}: Contiene únicamente estilos globales y compartidos por todas las páginas. Esto incluye la estructura de la cabecera, el pie de página, los botones base (\texttt{.cta-button}), los formularios (\texttt{.form-group}) y las tarjetas de contenido (\texttt{.sub-view-card}).
    \item \path{LandingPage.css}: Se creó un archivo de estilos específico para la página de inicio. Aquí se alojan los estilos de la sección principal (\textit{Hero}), como la cuadrícula de dos columnas, los títulos grandes, la imagen central y los efectos de animación del logotipo.
\end{itemize}

Esta organización garantiza que, si se modifica el tamaño del título en la página de inicio, esto no afecte a los títulos de la página de registro o del panel de control.

\section{Implementación de las pantallas principales}

\subsection{Página de inicio (Landing Page)}

La página de inicio es la carta de presentación del producto. Se diseñó con un enfoque \textit{responsive} (adaptable a móviles) y está compuesta por:

\begin{itemize}
    \item Un encabezado de navegación (\texttt{Header}) con el logotipo y botones para iniciar sesión o registrarse.
    \item Una sección central (\textit{Hero}) con el eslogan del proyecto: \textit{«Structure your lessons. See the progress.»}, donde la segunda frase se resalta utilizando el color secundario definido en la paleta.
    \item Una tarjeta de imagen con un efecto de sombra y esquinas redondeadas, que contiene la fotografía de un profesor o entorno educativo.
    \item Un pie de página (\texttt{Footer}) con los derechos de autor.
\end{itemize}

\subsection{Pantallas de autenticación y gestión de usuarios}

Se han desarrollado las siguientes vistas, todas ellas construidas sobre un componente de tarjeta (\texttt{sub-view-card}) que garantiza una apariencia consistente:

\begin{enumerate}
    \item \textbf{Registro (\texttt{RegisterPage}):} Un formulario dinámico que permite al usuario seleccionar su rol (Profesor, Alumno o Padre). Dependiendo de la selección, el formulario muestra u oculta campos específicos (por ejemplo, el teléfono para el profesor o el método de pago para el padre), tal como se definió en el contrato de la API.
    \item \textbf{Inicio de sesión (\texttt{LoginPage}):} Un formulario clásico de correo electrónico y contraseña, con un enlace para recuperar la contraseña en caso de olvido.
    \item \textbf{Recuperación de contraseña (\texttt{ForgotPasswordPage}):} Una vista que permite solicitar un enlace de restablecimiento. Al enviar el formulario, se muestra un mensaje de confirmación en la misma pantalla.
    \item \textbf{Verificación de correo (\texttt{EmailVerificationPage}):} Una pantalla informativa que se muestra tras un registro exitoso, indicando al usuario que revise su bandeja de entrada para activar su cuenta.
\end{enumerate}

\section{Resultados y siguientes pasos}

Al concluir esta fase de desarrollo del frontend, se han conseguido los siguientes hitos:

\begin{itemize}
    \item La interfaz web es completamente navegable gracias a React Router.
    \item El código está modularizado, lo que facilita la incorporación de nuevos desarrolladores (o la continuación del trabajo por parte de los mismos autores).
    \item Se ha implementado una identidad visual coherente con los diseños de Figma.
    \item Las pantallas de inicio de sesión y registro están listas para ser conectadas a la API de Spring Boot en el próximo \textit{sprint}.
\end{itemize}

El siguiente paso lógico es la creación de componentes de interfaz de usuario (UI Kit) reutilizables (botones, entradas de texto, tarjetas) en Storybook, vinculándolos a los diseños de Figma, y comenzar a construir el panel de control (\textit{Dashboard}) para el profesor, donde se gestionarán las clases y los alumnos.
\end{document}